\documentclass[10pt,a4paper]{article}
\usepackage[margin=2.2cm]{geometry}
\usepackage{fontspec}
\setmonofont{DejaVu Sans Mono}[Scale=0.84]
\newfontfamily\unifont{DejaVu Serif}[Scale=0.9]
\usepackage{amsmath,amssymb,mathtools}
\setcounter{MaxMatrixCols}{20}
\usepackage{longtable,booktabs,array}
\usepackage{fancyvrb}
\usepackage[dvipsnames]{xcolor}
\usepackage[colorlinks=true,linkcolor=NavyBlue,urlcolor=NavyBlue]{hyperref}
\usepackage{microtype}
\usepackage{needspace}
\usepackage{tocloft}\setlength{\cftbeforesecskip}{0.45em}
\setlength{\parindent}{0pt}\setlength{\parskip}{0.55em}
\setlength{\emergencystretch}{3em}
\renewcommand{\arraystretch}{1.15}
\allowdisplaybreaks
\begin{document}
\title{7+1: symbols and their definitions\\[0.4em]\large Every symbol used in the notebook 7+1.nb and in its documentation, with its defining equation}
\author{Repository: github.com/once-ere/7+1 — prepared for Nash\\ Notebook: 7+1.nb (Wolfram Language 15.0.1)\\ Generated by tools/make\_docs.py on 2026-10-26}
\date{}
\maketitle
\begin{abstract}
This is a self-contained dictionary. For each symbol it gives the mathematical notation, the Mathematica name in 7+1.nb, the defining equation and its meaning, and, where it exists, the name used in the earlier notebook 4+4.nb.
\end{abstract}
\tableofcontents
\Needspace{12\baselineskip}\section*{1. Conventions used in every definition below}\addcontentsline{toc}{section}{1. Conventions used in every definition below}
\begin{itemize}
\item Indices $A,B,K,L\in\lbrace0,1,\dots,7\rbrace$ label the coordinates $x_A$, the octonion units $e_A$ and the Clifford generators; $i,j\in\lbrace1,\dots,7\rbrace$; $h\in\lbrace1,2,3\rbrace$; $k\in\lbrace4,\dots,7\rbrace$ where stated.
\item Mathematica lists are 1-based: the entry $\eta_{AB}$ of a metric is \texttt{η[[A+1, B+1]]}.
\item $1_n$ is the $n\times n$ identity matrix (\texttt{IdentityMatrix[n]}, named \texttt{ID4}, \texttt{ID8}, \texttt{ID16}, \texttt{ID32}); $0$ in a block matrix is a zero block; $M^{\mathsf T}$ is the transpose, $\overline M$ the complex conjugate, $M^\dagger$ the conjugate transpose; $\delta_{AB}$ is the Kronecker delta; $\lbrace X,Y\rbrace=XY+YX$, $[X,Y]=XY-YX$.
\item $\left(\begin{smallmatrix}a&b\\c&d\end{smallmatrix}\right)$ is a block matrix (Mathematica \texttt{ArrayFlatten[\{\{a, b\}, \{c, d\}\}]}); $M\otimes N$ is the Kronecker product (\texttt{KroneckerProduct[M, N]}).
\item An index 8 on a family of eight Clifford generators $G_0,\dots,G_7$ means the ordered product $G_8=G_0G_1\cdots G_7$ (the convention of 4+4.nb).
\item $Cl(p,q)$: $p$ generators square to $+1$ and $q$ to $-1$ ($v^2=Q(v)$, $Q$ of signature $(p,q)$).
\end{itemize}
\Needspace{12\baselineskip}\section*{2. Coordinates and metrics}\addcontentsline{toc}{section}{2. Coordinates and metrics}
{\small
\begin{longtable}{>{\raggedright\arraybackslash}p{\dimexpr 0.1277\linewidth-2\tabcolsep\relax}>{\raggedright\arraybackslash}p{\dimexpr 0.1809\linewidth-2\tabcolsep\relax}>{\raggedright\arraybackslash}p{\dimexpr 0.3511\linewidth-2\tabcolsep\relax}>{\raggedright\arraybackslash}p{\dimexpr 0.3404\linewidth-2\tabcolsep\relax}}
\toprule
\textbf{symbol} & \textbf{Mathematica} & \textbf{definition} & \textbf{meaning} \\
\midrule\endfirsthead
\toprule
\textbf{symbol} & \textbf{Mathematica} & \textbf{definition} & \textbf{meaning} \\
\midrule\endhead
\endfoot
\bottomrule\endlastfoot
$X$ & \texttt{X} & $X=(x_0,x_1,\dots,x_7)$ & coordinates; also the components of a symbolic octonion or vector \\
$x_0$ & \texttt{x0} & — & 7+1.nb: ordinary time, octonion unit $e_0$. 4+4.nb: hidden space-like dimension \\
$x_1,x_2,x_3$ & \texttt{x1, x2, x3} & — & ordinary 3-space (both notebooks); octonion units $e_1,e_2,e_3$ \\
$x_4$ & \texttt{x4} & — & 7+1.nb: auxiliary space-like, $e_4=\ell$. 4+4.nb: ordinary time \\
$x_5,x_6,x_7$ & \texttt{x5, x6, x7} & — & 7+1.nb: auxiliary space-like, $e_5,e_6,e_7$. 4+4.nb: auxiliary time-like ("superluminal deflating time") \\
$\eta^{(4,4)}$ & \texttt{η4488} & $\mathrm{diag}(1,1,1,1,-1,-1,-1,-1)$ & metric of 4+4.nb (space-like $+1$, time-like $-1$); split-octonion norm \\
$\eta^{(7,1)}$ & \texttt{η71} & $\mathrm{diag}(-1,1,1,1,1,1,1,1)$ & 7+1 metric in the sign convention of 4+4.nb; signature (7,1) \\
$\eta^{(1,7)}$ & \texttt{η17} & $-\eta^{(7,1)}=\mathrm{diag}(1,-1,-1,-1,-1,-1,-1,-1)$ & 7+1 metric, octonionic convention: $x^{\mathsf T}\eta^{(1,7)}x=\mathrm{Re}(xx)$ \\
$\eta_{(1,7)}^{AB}$ & \texttt{η17inv} & $(\eta^{(1,7)})^{-1}=\eta^{(1,7)}$ & inverse metric, raises indices \\
$g^{O}$ & \texttt{gO} & $1_8$ & norm metric of the octonions, $N(x)=x^{\mathsf T}g^{O}x$ \\
$g$ & \texttt{gAB44} & $\eta^{(4,4)}$ & norm metric of the split octonions (4+4.nb: \texttt{g\_AB}) \\
$\eta^{(5,3)}$ & \texttt{η53} & $\mathrm{diag}(1,-1,-1,-1,1,1,1,1)$ & $\mathrm{Re}(xx)$ of the split octonions \\*
\end{longtable}}
\Needspace{12\baselineskip}\section*{3. Objects reproduced from 4+4.nb}\addcontentsline{toc}{section}{3. Objects reproduced from 4+4.nb}
{\footnotesize
\begin{longtable}{>{\raggedright\arraybackslash}p{\dimexpr 0.1277\linewidth-2\tabcolsep\relax}>{\raggedright\arraybackslash}p{\dimexpr 0.1915\linewidth-2\tabcolsep\relax}>{\raggedright\arraybackslash}p{\dimexpr 0.3830\linewidth-2\tabcolsep\relax}>{\raggedright\arraybackslash}p{\dimexpr 0.2979\linewidth-2\tabcolsep\relax}}
\toprule
\textbf{symbol} & \textbf{Mathematica} & \textbf{definition} & \textbf{meaning (4+4.nb name)} \\
\midrule\endfirsthead
\toprule
\textbf{symbol} & \textbf{Mathematica} & \textbf{definition} & \textbf{meaning (4+4.nb name)} \\
\midrule\endhead
\endfoot
\bottomrule\endlastfoot
$\epsilon_{hpq4}$ & \texttt{Qa[h,p,q]} & \texttt{Signature[\{h, p, q, 4\}]} & sign of the permutation (Levi-Civita symbol) \\
 & \texttt{Qb[h,p,q]} & $\delta_{p4}\delta_{qh}-\delta_{ph}\delta_{q4}$ & auxiliary symbol \\
$(s_h)_{pq}$ & \texttt{s4by4[h]} & $\epsilon_{hpq4}-(\delta_{p4}\delta_{qh}-\delta_{ph}\delta_{q4})$ & self-dual antisymmetric 4×4 (\texttt{SelfDualAntiSymmetric}) \\
$(t_h)_{pq}$ & \texttt{t4by4[h]} & $\epsilon_{hpq4}+(\delta_{p4}\delta_{qh}-\delta_{ph}\delta_{q4})$ & anti-self-dual antisymmetric 4×4 (\texttt{AntiSelfDualAntiSymmetric}) \\
$\sigma$ & \texttt{σ} & $\left(\begin{smallmatrix}0&1_4\\1_4&0\end{smallmatrix}\right)$ & split metric on spinor space (\texttt{σ}, \texttt{σ88}) \\
$\zeta_A$ & \texttt{ζ[A]} & $\zeta_0=1_8$; $\zeta_h=\left(\begin{smallmatrix}0&s_h\\s_h&0\end{smallmatrix}\right)$; $\zeta_{7-h}=\left(\begin{smallmatrix}0&t_h\\-t_h&0\end{smallmatrix}\right)$; $\zeta_7=\zeta_1\cdots\zeta_6$ & 8×8 building blocks (\texttt{ζ[A]}) \\
$\tau_A$ & \texttt{τ[A]} & $\tau_A=\zeta_A$ & split-octonion matrices (\texttt{τ[A]}) \\
$\bar\tau_A$ & \texttt{τbar[A]} & $\bar\tau_0=1_8$, $\bar\tau_A=\sigma\tau_A^{\mathsf T}\sigma$ & conjugate matrices (symbolized tau-bar) \\
$T_A$ & \texttt{T16[A]} & $\left(\begin{smallmatrix}0&\bar\tau_A\\\tau_A&0\end{smallmatrix}\right)$, $A=0,\dots,7$ & 16×16 Dirac matrices, $\lbrace T_A,T_B\rbrace=2\eta^{(4,4)}_{AB}1_{16}$ (\texttt{T16\textasciicircum{}A[A]}) \\
$T_8$ & \texttt{T16[8]} & $T_0T_1\cdots T_7=\mathrm{diag}(-1_8,1_8)$ & chirality matrix (\texttt{T16\textasciicircum{}A[8]}) \\
$\sigma_{16}$ & \texttt{σ16} & $T_0T_1T_2T_3=\mathrm{diag}(-\sigma,\sigma)$ & (\texttt{σ16}) \\
$P_L,P_R$ & \texttt{PL44}, \texttt{PR44} & $\tfrac12(1_{16}-T_8)$, $\tfrac12(1_{16}+T_8)$ & chiral projectors (\texttt{P\_L}, \texttt{P\_R}) \\
$U_h$ & \texttt{uBasis[[h]]} & eigenvectors of $\sigma$ divided by $\sqrt2$ (order of \texttt{Eigensystem[σ]}) & orthonormal eigenbasis of $\sigma$ (4+4.nb: double-struck u) \\
 & \texttt{evalues}, \texttt{evecs} & \texttt{Eigensystem[σ]} & eigenvalues $(-1,-1,-1,-1,1,1,1,1)$ and eigenvectors of $\sigma$ \\
$u$ & \texttt{unit} & $U_8=(1,0,0,0,1,0,0,0)^{\mathsf T}/\sqrt2$, \texttt{hUSE = 8} & unit spinor: $\sigma u=u$, $u^{\mathsf T}u=u^{\mathsf T}\sigma u=1$ (\texttt{unit}, \texttt{hUSE}) \\
$E_{44}(h)$ & \texttt{E44[h]} & $(E_{44})_{Ba}=(U_h^{\mathsf T}\sigma\bar\tau_B)_a$ & vielbein (\texttt{E\textasciicircum{}A\_a[h]}) \\
$E_{44}(h)^{-1}$ & \texttt{E44inv[h]} & $(U_h^{\mathsf T}\sigma U_h)\,(\eta^{(4,4)}_{00}\tau_0U_h,\dots,\eta^{(4,4)}_{77}\tau_7U_h)$ (columns) & inverse vielbein (\texttt{E\textasciicircum{}a\_A[h]}) \\
$\varepsilon_B$ & \texttt{eA[B]}, \texttt{EA} & $B=1,\dots,8$, \texttt{EA = Array[eA, 8]} & split-octonion units; $\varepsilon_B\leftrightarrow\tau_{B-1}\leftrightarrow x_{B-1}$ \\
 & \texttt{gridx44} & \texttt{Grid} of $\varepsilon_B\varepsilon_C=\sum_D(g_{BB}E_{44}\tau_{B-1}E_{44}^{-1})_{DC}\varepsilon_D$ & split-octonion multiplication table (\texttt{gridx}) \\
 & \texttt{gridx44From4p4nb} & the table printed by 4+4.nb, Out[325] & reference for the comparison check \\
$L^{S}_A$ & \texttt{LS[A]} & $g_{AA}\,E_{44}(8)\,\tau_A\,E_{44}(8)^{-1}$ & left multiplication by $\varepsilon_{A+1}$ \\
$x\star_S y$ & \texttt{sMul[x, y]} & $\big(\sum_Ax_AL^{S}_A\big)y$ & split-octonion product of coefficient lists \\
$\bar x$, $N_S(x)$ & \texttt{sConj}, \texttt{sNorm2} & $(x_0,-x_1,\dots,-x_7)$; $x^{\mathsf T}\eta^{(4,4)}x$ & split conjugate and norm \\
 & \texttt{zeroDivisor} & $(\varepsilon_1+\varepsilon_5)(\varepsilon_1-\varepsilon_5)=0$ & a product of two nonzero split octonions \\*
\end{longtable}}
\Needspace{12\baselineskip}\section*{4. The octonions}\addcontentsline{toc}{section}{4. The octonions}
{\footnotesize
\begin{longtable}{>{\raggedright\arraybackslash}p{\dimexpr 0.1064\linewidth-2\tabcolsep\relax}>{\raggedright\arraybackslash}p{\dimexpr 0.2128\linewidth-2\tabcolsep\relax}>{\raggedright\arraybackslash}p{\dimexpr 0.4043\linewidth-2\tabcolsep\relax}>{\raggedright\arraybackslash}p{\dimexpr 0.2766\linewidth-2\tabcolsep\relax}}
\toprule
\textbf{symbol} & \textbf{Mathematica} & \textbf{definition} & \textbf{meaning} \\
\midrule\endfirsthead
\toprule
\textbf{symbol} & \textbf{Mathematica} & \textbf{definition} & \textbf{meaning} \\
\midrule\endhead
\endfoot
\bottomrule\endlastfoot
$\tau^{O}_A$ & \texttt{τO[A]} & $\tau^{O}_0=1_8$; $\tau^{O}_A=\tau_A$ ($A=1,2,3$); $\tau^{O}_A=\tau_A\sigma$ ($A=4,\dots,7$) & octonion 8×8 matrices \\
$\bar\tau^{O}_A$ & \texttt{τObar[A]} & $g^{O}(\tau^{O}_A)^{\mathsf T}g^{O}=(\tau^{O}_A)^{\mathsf T}$ & conjugate matrices; $\tau^{O}_A\bar\tau^{O}_B+\tau^{O}_B\bar\tau^{O}_A=2g^{O}_{AB}1_8$ \\
$E^{O}$ & \texttt{EO} & $(E^{O})_{Ba}=(u^{\mathsf T}g^{O}\bar\tau^{O}_B)_a$; $E^{O}=\eta^{(4,4)}E_{44}(8)$ & octonion vielbein (orthogonal) \\
$(E^{O})^{-1}$ & \texttt{EOinv} & $(E^{O})^{\mathsf T}=(u^{\mathsf T}g^{O}u)(g^{O}_{00}\tau^{O}_0u,\dots,g^{O}_{77}\tau^{O}_7u)$ & inverse vielbein \\
$e_A$ & \texttt{e[A]}, \texttt{eList} & $A=0,\dots,7$, \texttt{eList = Array[e, 8, 0]} & octonion units, $e_0=1$; $e_A\leftrightarrow\tau^{O}_A\leftrightarrow x_A$ \\
$L^{O}_A$ & \texttt{LO[A]} & $g^{O}_{AA}\,E^{O}\tau^{O}_A(E^{O})^{-1}$ & left multiplication by $e_A$: $e_Ae_K=\sum_L(L^{O}_A)_{LK}e_L$ \\
 & \texttt{gridO} & \texttt{Grid} of $e_Ae_K$, same formula as \texttt{gridx} & octonion multiplication table \\
 & \texttt{gridOExpected} & table computed by \texttt{tools/ref71.py} & reference for the comparison check \\
$xy$ & \texttt{oMul[x, y]} & $\big(\sum_Ax_AL^{O}_A\big)y$ & octonion product of coefficient lists \\
$\bar x$ & \texttt{oConj[x]} & $(x_0,-x_1,\dots,-x_7)$ & conjugate \\
$N(x)$ & \texttt{oNorm2[x]} & $x^{\mathsf T}g^{O}x=x_0^2+\dots+x_7^2$ & norm \\
$\mathrm{Re}\,x$ & \texttt{oRe[x]} & $x_0$ & real part \\
$x^{-1}$ & \texttt{oInv[x]} & $\bar x/N(x)$ & inverse \\
$[x,y,z]$ & \texttt{oAssoc[x, y, z]} & $(xy)z-x(yz)$ & associator \\
 & \texttt{toE[x]}, \texttt{ev[A]} & $\sum_Ax_Ae_A$; \texttt{UnitVector[8, A+1]} & display; coefficient list of $e_A$ \\
 & \texttt{xs}, \texttt{ys}, \texttt{zs} & \texttt{Array[a, 8, 0]}, \texttt{Array[b, 8, 0]}, \texttt{Array[c, 8, 0]} & general symbolic octonions \\
 & \texttt{fanoLines}, \texttt{orientedLines} & $\lbrace i,j,k\rbrace$ with $e_ie_j=\pm e_k$; $(a,b,c)$ with $e_ae_b=+e_c$ & the 7 Fano lines \\
 & \texttt{fanoPicture} & centre $e_1$; vertices $e_2,e_4,e_7$; midpoints $e_6,e_3,e_5$ & picture of the Fano plane \\
 & \texttt{nonAssociativeTriples} & basis triples with $[e_A,e_B,e_K]\neq0$ & there are 168 \\
 & \texttt{xNum}, \texttt{yNum}, \texttt{zNum} & $e_0+2e_1-e_4+3e_7$, $-2e_0+e_2+e_5$, $e_3+e_6$ & numerical example \\
$\ell$ & \texttt{ev[4]} & $\ell=e_4$ & Cayley–Dickson doubling unit, $\ell^2=-1$ \\
$pq$, $\bar p$ & \texttt{qMul}, \texttt{qConj} & product in $\mathrm{span}\lbrace e_0,\dots,e_3\rbrace$; $(p_0,-p_1,-p_2,-p_3)$ & quaternions ($e_1e_2=-e_3$) \\
$(a,b)$ & \texttt{asPair[x]} & $x=a+\ell b$, $a=(x_0,\dots,x_3)$, $b=(x_4,\dots,x_7)$ & octonion as a pair of quaternions \\
 & \texttt{cdMul[γ]} & $(a,b)(c,d)=(ac+\gamma\,d\bar b,\ \bar a d+cb)$ & Cayley–Dickson; $\gamma=-1$ octonions, $+1$ split \\
 & \texttt{cdPrompt}, \texttt{pm} & $(ac-\bar d b,\ \bar a d+cb)$ & the formula of [1A] (fails) \\
 & \texttt{promptX}, \texttt{promptY} & $e_0+e_4=(1,1)$, $e_1=(i,0)$ & its counterexample \\
$s_{AB}$ & \texttt{sAB[A, B]} & $-1$ if $A\ge4$ and $B\ge4$, else $+1$ & octonion product $=s_{AB}\times$ split product \\
$c,\ \varphi$ & \texttt{cc}, \texttt{φ} & $c=(1,1,1,1,i,i,i,i)$, $\varphi(x)_A=c_Ax_A$ & isomorphism split$\otimes\mathbb C\to\mathbb O\otimes\mathbb C$ \\*
\end{longtable}}
\Needspace{12\baselineskip}\section*{5. Clifford algebras}\addcontentsline{toc}{section}{5. Clifford algebras}
{\footnotesize
\begin{longtable}{>{\raggedright\arraybackslash}p{\dimexpr 0.1064\linewidth-2\tabcolsep\relax}>{\raggedright\arraybackslash}p{\dimexpr 0.2128\linewidth-2\tabcolsep\relax}>{\raggedright\arraybackslash}p{\dimexpr 0.4255\linewidth-2\tabcolsep\relax}>{\raggedright\arraybackslash}p{\dimexpr 0.2553\linewidth-2\tabcolsep\relax}}
\toprule
\textbf{symbol} & \textbf{Mathematica} & \textbf{definition} & \textbf{meaning} \\
\midrule\endfirsthead
\toprule
\textbf{symbol} & \textbf{Mathematica} & \textbf{definition} & \textbf{meaning} \\
\midrule\endhead
\endfoot
\bottomrule\endlastfoot
$T^{O}_A$ & \texttt{T16O[A]} & $\left(\begin{smallmatrix}0&\bar\tau^{O}_A\\\tau^{O}_A&0\end{smallmatrix}\right)$ & $Cl(8,0)$: $\lbrace T^{O}_A,T^{O}_B\rbrace=2\delta_{AB}1_{16}$ \\
$T^{O}_8$ & \texttt{T16O[8]} & $T^{O}_0\cdots T^{O}_7=T_8=\mathrm{diag}(-1_8,1_8)$ & volume element of $Cl(8,0)$ \\
$S_{44}$ & \texttt{S44} & $T_4T_5T_6T_7=\mathrm{diag}(\sigma,\sigma)$ & $T^{O}_k=T_kS_{44}$ for $k\ge4$ \\
$\Gamma^{(1,7)}_A$ & \texttt{Γ17[A]} & $\Gamma^{(1,7)}_0=T^{O}_0=T_0$; $\Gamma^{(1,7)}_i=T^{O}_0T^{O}_i$ ($=T_0T_i$, $i\le3$; $=T_0T_iS_{44}$, $i\ge4$) & REAL 7+1 generators: $\lbrace\Gamma^{(1,7)}_A,\Gamma^{(1,7)}_B\rbrace=2\eta^{(1,7)}_{AB}1_{16}$ \\
$\Gamma^{(1,7)}_8$ & \texttt{Γ17[8]} & $\Gamma^{(1,7)}_0\cdots\Gamma^{(1,7)}_7=\left(\begin{smallmatrix}0&-1_8\\1_8&0\end{smallmatrix}\right)$ & volume element, square $-1_{16}$ \\
 & \texttt{products256} & $\Gamma^{(1,7)}_{A_1}\cdots\Gamma^{(1,7)}_{A_k}$, $A_1<\dots<A_k$ & the 256 basis products \\
$\tau^{C}_A$, $\bar\tau^{C}_A$ & \texttt{τC[A]}, \texttt{τCbar[A]} & $\tau^{C}_0=\bar\tau^{C}_0=i1_8$; $\tau^{C}_k=\tau^{O}_k$, $\bar\tau^{C}_k=\bar\tau^{O}_k$ & complexified 8×8 matrices \\
$\Gamma^{(7,1)}_A$ & \texttt{Γ71[A]} & $\left(\begin{smallmatrix}0&\bar\tau^{C}_A\\\tau^{C}_A&0\end{smallmatrix}\right)$; $\Gamma^{(7,1)}_0=iT_0$, $\Gamma^{(7,1)}_k=T^{O}_k$ & COMPLEX generators: $\lbrace\Gamma^{(7,1)}_A,\Gamma^{(7,1)}_B\rbrace=2\eta^{(7,1)}_{AB}1_{16}$ \\
$\Gamma^{(7,1)}_8$ & \texttt{Γ71[8]} & $\Gamma^{(7,1)}_0\cdots\Gamma^{(7,1)}_7=iT_8$ & volume element, square $-1_{16}$ \\
$S$ & \texttt{Ssim} & $(1_{16}-i\Gamma^{(1,7)}_0)/\sqrt2$ & unitary; $\Gamma^{(7,1)}_A=S(i\Gamma^{(1,7)}_A)S^{-1}$ \\
 & \texttt{intertwiners[gs]} & \texttt{NullSpace} basis of $\lbrace M:M\overline g=gM\ \forall g\rbrace$ & Schur intertwiners \\
 & \texttt{indicator[gs]} & sign of the scalar $M\overline M$ & $+1$: $\mathbb R(16)$; $-1$: $\mathbb H(8)$ \\
$M_0$ & \texttt{M71} & intertwiners of $\Gamma^{(7,1)}$, spanned by $\Gamma^{(1,7)}_8$ & $M_0\overline{M_0}=-1_{16}$ \\
$\varepsilon$ & \texttt{ε2} & $\left(\begin{smallmatrix}0&1\\-1&0\end{smallmatrix}\right)$ & $\varepsilon^2=-1_2$ \\
$\widehat\Gamma^{(7,1)}_A$ & \texttt{Γ71R32[A]} & $\varepsilon\otimes\Gamma^{(1,7)}_A$ & REAL 32×32 generators with $\eta^{(7,1)}$ \\
$Q_i$, $Q_j$, $Q_k$ & \texttt{Qi}, \texttt{Qj}, \texttt{Qk} & $\varepsilon\otimes1_{16}$, $\sigma_1\otimes\Gamma^{(1,7)}_8$, $\sigma_3\otimes\Gamma^{(1,7)}_8$ & quaternion units commuting with all $\widehat\Gamma^{(7,1)}_A$ \\
 & \texttt{commutantBasis}, \texttt{commutantTime} & \texttt{NullSpace} of $\lbrace X: X\widehat\Gamma^{(7,1)}_A=\widehat\Gamma^{(7,1)}_AX\rbrace$; its timing & dimension 4: $\mathbb H$ \\
$\chi^{(7,1)}$ & \texttt{χ71} & $-i\Gamma^{(7,1)}_8=T_8$ & chirality matrix ($\Gamma^{(7,1)}$ basis) \\
 & \texttt{PL71}, \texttt{PR71} & $\tfrac12(1_{16}\mp\chi^{(7,1)})=P_L,P_R$ & 7+1 chiral projectors \\
$\chi^{(1,7)}$ & \texttt{χ17}, \texttt{PL17}, \texttt{PR17} & $i\Gamma^{(1,7)}_8$; $\tfrac12(1_{16}\mp\chi^{(1,7)})$ & complex chirality and projectors (rank 8) \\
$E^{O}_{16}$ & \texttt{EO16} & $\mathrm{diag}(E^{O},E^{O})$ & rotation to the octonion basis \\
$\Gamma^{e}_A$ & \texttt{Γ17e[A]} & $E^{O}_{16}\Gamma^{(1,7)}_A(E^{O}_{16})^{\mathsf T}$; $\Gamma^{e}_0=\left(\begin{smallmatrix}0&1_8\\1_8&0\end{smallmatrix}\right)$, $\Gamma^{e}_i=\mathrm{diag}(L^{O}_i,-L^{O}_i)$ & $\Gamma^{(1,7)}$ in the octonion basis \\
$\vec x$, $\psi$ & \texttt{xImag}, \texttt{ψs} & $(0,x_1,\dots,x_7)$; \texttt{Array[ψ, 8, 0]} & imaginary part; a general octonion \\*
\end{longtable}}
\Needspace{12\baselineskip}\section*{6. Examples and bookkeeping}\addcontentsline{toc}{section}{6. Examples and bookkeeping}
{\footnotesize
\begin{longtable}{>{\raggedright\arraybackslash}p{\dimexpr 0.1064\linewidth-2\tabcolsep\relax}>{\raggedright\arraybackslash}p{\dimexpr 0.2340\linewidth-2\tabcolsep\relax}>{\raggedright\arraybackslash}p{\dimexpr 0.4043\linewidth-2\tabcolsep\relax}>{\raggedright\arraybackslash}p{\dimexpr 0.2553\linewidth-2\tabcolsep\relax}}
\toprule
\textbf{symbol} & \textbf{Mathematica} & \textbf{definition} & \textbf{meaning} \\
\midrule\endfirsthead
\toprule
\textbf{symbol} & \textbf{Mathematica} & \textbf{definition} & \textbf{meaning} \\
\midrule\endhead
\endfoot
\bottomrule\endlastfoot
$\Sigma_{AB}$ & \texttt{Σ17[A, B]}, \texttt{pairs} & $\tfrac14[\Gamma^{(1,7)}_A,\Gamma^{(1,7)}_B]$; \texttt{pairs} lists the 28 pairs $A<B$ & generators of Spin(1,7) \\
$S(\beta)^{\pm1}$ & \texttt{Sβ}, \texttt{Sβinv} & $\cosh\tfrac\beta2\,1_{16}\pm\sinh\tfrac\beta2\,\Gamma^{(1,7)}_0\Gamma^{(1,7)}_1$ & boost along $x_1$ and its inverse \\
$R(\theta)^{\pm1}$ & \texttt{Rθ}, \texttt{Rθinv} & $\cos\tfrac\theta2\,1_{16}\pm\sin\tfrac\theta2\,\Gamma^{(1,7)}_4\Gamma^{(1,7)}_5$ & rotation in the $x_4$–$x_5$ plane and its inverse \\
$\Lambda$ & \texttt{Λboost}, \texttt{Λrot} & $\Lambda_{BK}=\tfrac1{16}\mathrm{tr}(\eta^{(1,7)}_{BB}\Gamma^{(1,7)}_BS\Gamma^{(1,7)}_KS^{-1})$ & Lorentz matrices of the two examples \\
$\Gamma^{A}$ & \texttt{Γ17up[A]} & $\eta_{(1,7)}^{AA}\Gamma^{(1,7)}_A$ & raised index \\
$\Psi$ & \texttt{Ψfun} & $(f_1(X),\dots,f_{16}(X))$, \texttt{f[k] @@ X} & 16 arbitrary functions of $x_0,\dots,x_7$ \\
$\Gamma^{A}\partial_A$ & \texttt{DiracOp17[ψ]} & $\sum_A\Gamma^{A}\,\partial\psi/\partial x_A$ & 7+1 Dirac operator \\
$\square$ & \texttt{waveOp17[ψ]} & $\sum_A\eta_{(1,7)}^{AA}\,\partial^2\psi/\partial x_A^2$ & 7+1 wave operator \\
$p,\ m$ & \texttt{pvec}, \texttt{mass} & $p=(3,2,2,0,0,0,0,0)$ (contravariant $p^A$), $m=1$ & $p^{\mathsf T}\eta^{(1,7)}p=m^2$ \\
$\not p$ & \texttt{pslash} & $\sum_Ap^A\Gamma^{(1,7)}_A$ & $\not p^{\,2}=m^21_{16}$ \\
$W$ & \texttt{realAmplitudes} & \texttt{NullSpace[pslash - mass ID16]} & eight real spinor amplitudes, $(\not p-m)w=0$ \\
$p\cdot x$ & \texttt{pvec.η17.X} & $p^{\mathsf T}\eta^{(1,7)}x$ & phase of the plane waves \\
 & \texttt{Ψwave} & $w\,e^{-i\,p\cdot x}$ & solves $(i\Gamma^{A}\partial_A-m)\Psi=0$; complex \\
$J$ & \texttt{antiCommutant} & \texttt{NullSpace} of $J\Gamma^{(1,7)}_A+\Gamma^{(1,7)}_AJ=0$ ($A=0,\dots,7$) & spanned by $\Gamma^{(1,7)}_8$ alone, $J^2<0$ \\
$k,\ \not k$ & \texttt{kvec}, \texttt{kslash} & $k=(1,1,0,0,0,0,0,0)$, $\not k=\sum_Ak^A\Gamma^{(1,7)}_A$ & null momentum, $\not k^{\,2}=0$ \\
 & \texttt{realNull} & \texttt{NullSpace[kslash]} & 8 real solutions of $\not k\,w=0$ \\
 & \texttt{ΨmasslessReal} & $w\cos(k\cdot x)$ & REAL solution of $\Gamma^{A}\partial_A\Psi=0$ \\
$\eta_{(7,1)}^{AB}$ & \texttt{η71inv} & $(\eta^{(7,1)})^{-1}=\eta^{(7,1)}$ & raises indices for $\widehat\Gamma$ \\
$\widehat\Gamma^{A}\partial_A$ & \texttt{DiracOp71R32[ψ]} & $\sum_A\eta_{(7,1)}^{AA}\widehat\Gamma^{(7,1)}_A\,\partial\psi/\partial x_A$ & real 32-component Dirac operator \\
 & \texttt{Ψfun32} & $(g_1(X),\dots,g_{32}(X))$, \texttt{g32[k] @@ X} & 32 arbitrary functions \\
$P$ & \texttt{pslash32} & $\sum_Ap^A\widehat\Gamma^{(7,1)}_A$ & $P^2=-m^2\,1_{32}$ \\
$w$ & \texttt{w32} & \texttt{Array[w, 32]} & 32 arbitrary real constants \\
 & \texttt{Ψ32} & $w\cos(p\cdot x)+\tfrac1mPw\sin(p\cdot x)$ & REAL solution of $(\widehat\Gamma^{A}\partial_A-m)\Psi=0$ \\
$\Gamma^{(5,3)}_A$ & \texttt{Γ53[A]} & $\Gamma^{(5,3)}_0=T_0$, $\Gamma^{(5,3)}_i=T_0T_i$ & real $Cl(5,3)$ with $\eta^{(5,3)}$ \\
$\Gamma^{(0,8)}_A$ & \texttt{Γ08[A]} & $T^{O}_8T^{O}_A$ & real $Cl(0,8)$ \\
 & \texttt{gens[p]} & $T^{O}_A$ for $A<p$, $iT^{O}_A$ for $A\ge p$ & complex generators of signature $(p,8-p)$ \\
 & \texttt{classification}, \texttt{claimed1B}, \texttt{correct} & $(p,q,\text{indicator})$, $p=8,\dots,0$; claimed and computed classes & classification table \\
 & \texttt{checks}, \texttt{check[name, test]} & association of results; stores \texttt{TrueQ[test]} and returns \texttt{name -> result} & verification bookkeeping \\
 & \texttt{failedChecks}, \texttt{allChecksPassed} & names of checks that are not True; \texttt{AllTrue[Values[checks], TrueQ]} & must be \texttt{\{\}} and \texttt{True} \\*
\end{longtable}}
\Needspace{12\baselineskip}\section*{7. Explicit matrices}\addcontentsline{toc}{section}{7. Explicit matrices}
The 4×4 blocks, the seven split matrices $\tau_A$, and the four new octonion matrices ($\tau^{O}_A=\tau_A$ for $A\le3$):
\begin{equation*}
s_1=\begin{pmatrix}0 & 0 & 0 & 1 \\ 0 & 0 & 1 & 0 \\ 0 & -1 & 0 & 0 \\ -1 & 0 & 0 & 0\end{pmatrix},\ s_2=\begin{pmatrix}0 & 0 & -1 & 0 \\ 0 & 0 & 0 & 1 \\ 1 & 0 & 0 & 0 \\ 0 & -1 & 0 & 0\end{pmatrix},\ s_3=\begin{pmatrix}0 & 1 & 0 & 0 \\ -1 & 0 & 0 & 0 \\ 0 & 0 & 0 & 1 \\ 0 & 0 & -1 & 0\end{pmatrix},
\end{equation*}
\begin{equation*}
t_1=\begin{pmatrix}0 & 0 & 0 & -1 \\ 0 & 0 & 1 & 0 \\ 0 & -1 & 0 & 0 \\ 1 & 0 & 0 & 0\end{pmatrix},\ t_2=\begin{pmatrix}0 & 0 & -1 & 0 \\ 0 & 0 & 0 & -1 \\ 1 & 0 & 0 & 0 \\ 0 & 1 & 0 & 0\end{pmatrix},\ t_3=\begin{pmatrix}0 & 1 & 0 & 0 \\ -1 & 0 & 0 & 0 \\ 0 & 0 & 0 & -1 \\ 0 & 0 & 1 & 0\end{pmatrix},
\end{equation*}
\begin{equation*}
\tau_1=\begin{pmatrix}0 & 0 & 0 & 0 & 0 & 0 & 0 & 1 \\ 0 & 0 & 0 & 0 & 0 & 0 & 1 & 0 \\ 0 & 0 & 0 & 0 & 0 & -1 & 0 & 0 \\ 0 & 0 & 0 & 0 & -1 & 0 & 0 & 0 \\ 0 & 0 & 0 & 1 & 0 & 0 & 0 & 0 \\ 0 & 0 & 1 & 0 & 0 & 0 & 0 & 0 \\ 0 & -1 & 0 & 0 & 0 & 0 & 0 & 0 \\ -1 & 0 & 0 & 0 & 0 & 0 & 0 & 0\end{pmatrix},\quad \tau_2=\begin{pmatrix}0 & 0 & 0 & 0 & 0 & 0 & -1 & 0 \\ 0 & 0 & 0 & 0 & 0 & 0 & 0 & 1 \\ 0 & 0 & 0 & 0 & 1 & 0 & 0 & 0 \\ 0 & 0 & 0 & 0 & 0 & -1 & 0 & 0 \\ 0 & 0 & -1 & 0 & 0 & 0 & 0 & 0 \\ 0 & 0 & 0 & 1 & 0 & 0 & 0 & 0 \\ 1 & 0 & 0 & 0 & 0 & 0 & 0 & 0 \\ 0 & -1 & 0 & 0 & 0 & 0 & 0 & 0\end{pmatrix},
\end{equation*}
\begin{equation*}
\tau_3=\begin{pmatrix}0 & 0 & 0 & 0 & 0 & 1 & 0 & 0 \\ 0 & 0 & 0 & 0 & -1 & 0 & 0 & 0 \\ 0 & 0 & 0 & 0 & 0 & 0 & 0 & 1 \\ 0 & 0 & 0 & 0 & 0 & 0 & -1 & 0 \\ 0 & 1 & 0 & 0 & 0 & 0 & 0 & 0 \\ -1 & 0 & 0 & 0 & 0 & 0 & 0 & 0 \\ 0 & 0 & 0 & 1 & 0 & 0 & 0 & 0 \\ 0 & 0 & -1 & 0 & 0 & 0 & 0 & 0\end{pmatrix},\quad \tau_4=\begin{pmatrix}0 & 0 & 0 & 0 & 0 & 1 & 0 & 0 \\ 0 & 0 & 0 & 0 & -1 & 0 & 0 & 0 \\ 0 & 0 & 0 & 0 & 0 & 0 & 0 & -1 \\ 0 & 0 & 0 & 0 & 0 & 0 & 1 & 0 \\ 0 & -1 & 0 & 0 & 0 & 0 & 0 & 0 \\ 1 & 0 & 0 & 0 & 0 & 0 & 0 & 0 \\ 0 & 0 & 0 & 1 & 0 & 0 & 0 & 0 \\ 0 & 0 & -1 & 0 & 0 & 0 & 0 & 0\end{pmatrix},
\end{equation*}
\begin{equation*}
\tau_5=\begin{pmatrix}0 & 0 & 0 & 0 & 0 & 0 & -1 & 0 \\ 0 & 0 & 0 & 0 & 0 & 0 & 0 & -1 \\ 0 & 0 & 0 & 0 & 1 & 0 & 0 & 0 \\ 0 & 0 & 0 & 0 & 0 & 1 & 0 & 0 \\ 0 & 0 & 1 & 0 & 0 & 0 & 0 & 0 \\ 0 & 0 & 0 & 1 & 0 & 0 & 0 & 0 \\ -1 & 0 & 0 & 0 & 0 & 0 & 0 & 0 \\ 0 & -1 & 0 & 0 & 0 & 0 & 0 & 0\end{pmatrix},\quad \tau_6=\begin{pmatrix}0 & 0 & 0 & 0 & 0 & 0 & 0 & -1 \\ 0 & 0 & 0 & 0 & 0 & 0 & 1 & 0 \\ 0 & 0 & 0 & 0 & 0 & -1 & 0 & 0 \\ 0 & 0 & 0 & 0 & 1 & 0 & 0 & 0 \\ 0 & 0 & 0 & 1 & 0 & 0 & 0 & 0 \\ 0 & 0 & -1 & 0 & 0 & 0 & 0 & 0 \\ 0 & 1 & 0 & 0 & 0 & 0 & 0 & 0 \\ -1 & 0 & 0 & 0 & 0 & 0 & 0 & 0\end{pmatrix},
\end{equation*}
\begin{equation*}
\tau_7=\begin{pmatrix}-1 & 0 & 0 & 0 & 0 & 0 & 0 & 0 \\ 0 & -1 & 0 & 0 & 0 & 0 & 0 & 0 \\ 0 & 0 & -1 & 0 & 0 & 0 & 0 & 0 \\ 0 & 0 & 0 & -1 & 0 & 0 & 0 & 0 \\ 0 & 0 & 0 & 0 & 1 & 0 & 0 & 0 \\ 0 & 0 & 0 & 0 & 0 & 1 & 0 & 0 \\ 0 & 0 & 0 & 0 & 0 & 0 & 1 & 0 \\ 0 & 0 & 0 & 0 & 0 & 0 & 0 & 1\end{pmatrix},
\end{equation*}
\begin{equation*}
\tau^{O}_4=\begin{pmatrix}0 & 1 & 0 & 0 & 0 & 0 & 0 & 0 \\ -1 & 0 & 0 & 0 & 0 & 0 & 0 & 0 \\ 0 & 0 & 0 & -1 & 0 & 0 & 0 & 0 \\ 0 & 0 & 1 & 0 & 0 & 0 & 0 & 0 \\ 0 & 0 & 0 & 0 & 0 & -1 & 0 & 0 \\ 0 & 0 & 0 & 0 & 1 & 0 & 0 & 0 \\ 0 & 0 & 0 & 0 & 0 & 0 & 0 & 1 \\ 0 & 0 & 0 & 0 & 0 & 0 & -1 & 0\end{pmatrix},\quad \tau^{O}_5=\begin{pmatrix}0 & 0 & -1 & 0 & 0 & 0 & 0 & 0 \\ 0 & 0 & 0 & -1 & 0 & 0 & 0 & 0 \\ 1 & 0 & 0 & 0 & 0 & 0 & 0 & 0 \\ 0 & 1 & 0 & 0 & 0 & 0 & 0 & 0 \\ 0 & 0 & 0 & 0 & 0 & 0 & 1 & 0 \\ 0 & 0 & 0 & 0 & 0 & 0 & 0 & 1 \\ 0 & 0 & 0 & 0 & -1 & 0 & 0 & 0 \\ 0 & 0 & 0 & 0 & 0 & -1 & 0 & 0\end{pmatrix},
\end{equation*}
\begin{equation*}
\tau^{O}_6=\begin{pmatrix}0 & 0 & 0 & -1 & 0 & 0 & 0 & 0 \\ 0 & 0 & 1 & 0 & 0 & 0 & 0 & 0 \\ 0 & -1 & 0 & 0 & 0 & 0 & 0 & 0 \\ 1 & 0 & 0 & 0 & 0 & 0 & 0 & 0 \\ 0 & 0 & 0 & 0 & 0 & 0 & 0 & 1 \\ 0 & 0 & 0 & 0 & 0 & 0 & -1 & 0 \\ 0 & 0 & 0 & 0 & 0 & 1 & 0 & 0 \\ 0 & 0 & 0 & 0 & -1 & 0 & 0 & 0\end{pmatrix},\quad \tau^{O}_7=\begin{pmatrix}0 & 0 & 0 & 0 & -1 & 0 & 0 & 0 \\ 0 & 0 & 0 & 0 & 0 & -1 & 0 & 0 \\ 0 & 0 & 0 & 0 & 0 & 0 & -1 & 0 \\ 0 & 0 & 0 & 0 & 0 & 0 & 0 & -1 \\ 1 & 0 & 0 & 0 & 0 & 0 & 0 & 0 \\ 0 & 1 & 0 & 0 & 0 & 0 & 0 & 0 \\ 0 & 0 & 1 & 0 & 0 & 0 & 0 & 0 \\ 0 & 0 & 0 & 1 & 0 & 0 & 0 & 0\end{pmatrix}.
\end{equation*}
\Needspace{12\baselineskip}\section*{8. Old names (4+4.nb) and new names (7+1.nb)}\addcontentsline{toc}{section}{8. Old names (4+4.nb) and new names (7+1.nb)}
{\small
\begin{longtable}{>{\raggedright\arraybackslash}p{\dimexpr 0.3191\linewidth-2\tabcolsep\relax}>{\raggedright\arraybackslash}p{\dimexpr 0.3191\linewidth-2\tabcolsep\relax}>{\raggedright\arraybackslash}p{\dimexpr 0.3617\linewidth-2\tabcolsep\relax}}
\toprule
\textbf{4+4.nb} & \textbf{7+1.nb} & \textbf{reason} \\
\midrule\endfirsthead
\toprule
\textbf{4+4.nb} & \textbf{7+1.nb} & \textbf{reason} \\
\midrule\endhead
\endfoot
\bottomrule\endlastfoot
\texttt{σ}, \texttt{σ88} & \texttt{σ} & unchanged \\
\texttt{τ[A]}, \texttt{ζ[A]}, \texttt{s4by4[h]}, \texttt{t4by4[h]} & unchanged & — \\
symbolized tau-bar & \texttt{τbar[A]} & 2-D symbol made with the Notation package; a plain name needs no package \\
\texttt{T16\textasciicircum{}A[A1]} (symbolized) & \texttt{T16[A]} & the superscript looked like a power and the letter A was not an index \\
\texttt{P\_L}, \texttt{P\_R} & \texttt{PL44}, \texttt{PR44} & the 7+1 projectors get their own names \texttt{PL71}, \texttt{PR71}, \texttt{PL17}, \texttt{PR17} \\
double-struck u & \texttt{uBasis} & plain name \\
\texttt{g\_AB} & \texttt{gAB44} & subscripted name; plain name \\
\texttt{E\textasciicircum{}A\_a[h]}, \texttt{E\textasciicircum{}a\_A[h]} (symbolized) & \texttt{E44[h]}, \texttt{E44inv[h]} & the super- and subscripts were part of the name, not indices \\
\texttt{eA[1..8]} (1-based) & \texttt{e[0..7]} (0-based) for the octonions & aligns $e_A$ with $\tau^{O}_A$ and $x_A$: no $B-1$ shift \\
\texttt{gridx} & \texttt{gridx44}, and the new \texttt{gridO} & two tables now \\*
\end{longtable}}
\end{document}
